\subsection{\texorpdfstring{正则代数 \(^{1}\) 或正规代数中基域的可分扩张}{正则代数或正规代数中基域的可分扩张（1）}}
\footnotetext[1]{在本节中，域 \(k\) 上的一个代数称为正则的，如果它是正则环。特别地，\(k\) 的每个扩张都是正则代数。这一概念不应与 A, V, 第 135 页, 定义 2 中引入的正则扩张的概念相混淆，后者在此不会用到。}

引理 1.--- 设 A 为 Noether 环，是 Noether 子环的一个递增滤过族 \((\mathrm{A}_{\alpha})_{\alpha\in\mathbb{I}}\) 的并。

\begin{enumerate}
\def\labelenumi{\alph{enumi})}
\item
  若环 \(\Lambda_{\alpha}\) 都正则，且 A 是平坦 \(\mathrm{A}_{\alpha}\) -模（对所有 \(\alpha \in \mathrm{I}\)），则环 A 正则。
\item
  若环 \(\mathrm{A}_{\alpha}\) 都是正规的 (§ 1, 第 8 节)，则 A 是正规的。
\end{enumerate}

\begin{enumerate}
\def\labelenumi{\alph{enumi})}
\item
  设 \(\mathfrak{m}\) 为 A 的一个极大理想；对每个 \(\alpha \in I\)，用 \(\mathfrak{m}_{\alpha}\) 表示理想 \(\mathfrak{m} \cap A_{\alpha}\)（\(A_{\alpha}\) 中的）。由于 A 是 Noether 的，\(\mathfrak{m}\) 是有限型的，且存在 I 的一个元素 \(\alpha\) 使 \(\mathfrak{m} = A\mathfrak{m}_{\alpha}\)，于是 A-模 \((A_{\alpha}/\mathfrak{m}_{\alpha}) \otimes_{A_{\alpha}} A\) 同构于 \(A/m\)。由于 A 在 \(A_{\alpha}\) 上平坦，有 \(\mathrm{dp}_{A}(A/m) \leqslant \mathrm{dp}_{A_{\alpha}}(A_{\alpha}/\mathfrak{m}_{\alpha})\) (A, X, 第 141 页, lcmc 2)。由于环 \(A_{\alpha}\) 正则，于是由 \(\S 4\) (第 2 节, 命题 4) 得出 A 正则。
\item
  由于 \(A_{\alpha}\) 都是约化的，A 是约化的。设 a 和 b 为 A 的元素，使 b 不是零因子，且 a/b（A 的全分式环中的元素）在 \(\Lambda\) 上是整的。存在一个首一多项式 \(P \in A[X]\) 使 \(P(a/b) = 0\) 。设 \(\alpha\) 为 I 的一个元素，使环 \(A_{\alpha}\) 包含 a、b 及 P 的系数。由于 \(A_{\alpha}\) 正规，存在 \(c \in A_{\alpha}\) 使 a = bc。于是 a/b = c 属于 \(\Lambda\)，从而 \(\Lambda\) 是正规的。
\end{enumerate}

引理 2. -- 设 \(k\) 为域，K 和 L 为 \(k\) 的扩张。设 K 是有限型的，且扩张 K 与 L 之一是可分的。则环 \(\mathrm{K} \otimes_{k} \mathrm{L}\) 是正则的。

用 A 表示环 K ⊗k L；它由第 2 节的命题 2 的推论是 Noether 的。先设扩张 K 是可分的。由 A, V, 第 130 页的推论，存在 K 的一个超越基 t = (t₁, \ldots..., tₙ)，使 K 是纯扩张 k(t) 的有限可分扩张。环 E = k(t) ⊗k L 同构于 L 上一个多项式环的分式环，是正则的 (§ 4, 第 2 节, 命题 4 的推论 5)；对 E 的每个素理想 p，环 κ(p) ⊗E A 同构于 κ(p) ⊗k(t) K，从而同构于域的有限积 (A, V, 第 35 页, 定义 1, 第 34 页, 定理 4, 及第 33 页, 命题 5)。由于 A 是自由 E-模，它由 § 4, 第 5 节的命题 9 的推论是正则环。

现在设 L 在 k 上可分。由证明的第一部分，环 \(K \otimes_{k} L'\) 对 L 的每个有限生成子扩张 \(L'\) 都正则。我们应用引理 1 得出结论。

命题 5.--- 设 k 为域，A 为 k-代数，K 为 k 的扩张。假设 A 是本质有限型的，或扩张 K 是有限型的。

\begin{enumerate}
\def\labelenumi{\alph{enumi})}
\item
  若环 A \(_{(K)}\) 正则（相应地正规），则环 A 正则（相应地正规）。
\item
  若环 \(\Lambda\) 正则（相应地正规）且 K 在 \(k\) 上的扩张可分，则环 \(\mathrm{A}_{(\mathrm{K})}\) 正则（相应地正规）。
\end{enumerate}

A-模 \(A_{(K)}\) 是自由的，从而是忠实平坦的；于是断言 a) 由 I, § 3, 第 5 节的命题 8 的推论及 § 4, 第 5 节的命题 8 得出（相应地由 § 1, 第 10 节的定理 4 的推论 2 得出）。在 b) 的假设下，对 A 的每个素理想 p，环 \(\kappa(\mathfrak{p})\otimes_{k}K\) 是正则的 (引理 2)，且更加是正规的 (§ 4, 第 2 节, 命题 4 的推论 2)。于是断言 b) 由 § 4, 第 5 节的命题 9 的推论得出（相应地由 § 1, 第 10 节的定理 4 的推论 3 得出）。
% semantic-fragments: legacy-input-seam
\relax{}
